% ==============================================================================
%  PLAN DE ENTREGAS — ÓPERA PRIMA
%  Versión 2.0 — 4 semanas · 20 días hábiles
%  Sprint 1: MVP (15 días) · Sprint 2: Correcciones (15 días)
%  Equipo: Santiago Montaño (Frontend) · Camilo Sotelo (Backend)
%  Coordinación: Alexander Paz
%  Compila con: pdflatex (2 pasadas)
%  Diseño editorial con identidad visual de Ópera Prima
%  Paleta: Rosa coral #F65B7F · Verde selva #1A4A3C · Lavanda #5E3A8A
% ==============================================================================

\documentclass[11pt]{article}

% ── Encoding & Language ──
\usepackage[utf8]{inputenc}
\usepackage[spanish]{babel}
\usepackage[T1]{fontenc}

% ── Typography ──
\usepackage[scaled]{helvet}
\renewcommand{\familydefault}{\sfdefault}

% ── Geometry ──
\usepackage[
  a4paper,
  margin=2.5cm,
  bottom=3cm,
  headheight=2cm,
  footskip=1.5cm
]{geometry}

% ── Colors (Opera Prima) ──
\usepackage{xcolor}
\definecolor{op-rosa}{HTML}{F65B7F}
\definecolor{op-rosa-dark}{HTML}{D94268}
\definecolor{op-verde}{HTML}{1A4A3C}
\definecolor{op-lavanda}{HTML}{5E3A8A}
\definecolor{op-nearblack}{HTML}{111111}
\definecolor{op-surface}{HTML}{FAFAF9}
\definecolor{op-white}{HTML}{FFFFFF}
\definecolor{op-text}{HTML}{111111}
\definecolor{op-textsec}{HTML}{52525B}
\definecolor{op-textmuted}{HTML}{A1A1AA}

% ── Graphics ──
\usepackage{graphicx}
\usepackage{tikz}
\usetikzlibrary{shapes.geometric, arrows.meta, positioning}
\usepackage{booktabs}
\usepackage{array}
\usepackage{colortbl}
\usepackage{multirow}

% ── Boxes ──
\usepackage{tcolorbox}
\tcbuselibrary{skins, breakable}

\tcbset{
  op-box/.style={
    colback=op-surface,
    colframe=op-nearblack,
    boxrule=0.4mm,
    arc=0mm,
    outer arc=0mm,
    left=5mm,
    right=5mm,
    top=3mm,
    bottom=3mm,
    sharp corners,
  },
  op-fase/.style={
    colback=op-nearblack,
    colframe=op-nearblack,
    coltext=op-white,
    boxrule=0mm,
    arc=0mm,
    outer arc=0mm,
    left=5mm,
    right=5mm,
    top=2.5mm,
    bottom=2.5mm,
    sharp corners,
    fontupper=\Large\bfseries,
  },
  op-aviso/.style={
    colback=op-rosa!12,
    colframe=op-rosa,
    boxrule=0.3mm,
    arc=0mm,
    outer arc=0mm,
    left=4mm,
    right=4mm,
    top=2mm,
    bottom=2mm,
    sharp corners,
  },
  op-mvp/.style={
    colback=op-verde!10,
    colframe=op-verde,
    boxrule=0.3mm,
    arc=0mm,
    outer arc=0mm,
    left=4mm,
    right=4mm,
    top=2mm,
    bottom=2mm,
    sharp corners,
  },
}

% ── Section Headings ──
\usepackage{titlesec}
\titleformat{\section}
  {\LARGE\bfseries\color{op-nearblack}}
  {\thesection}{1em}{}[\vspace{-0.3em}{\color{op-rosa}\rule{\textwidth}{0.4mm}}]
\titleformat{\subsection}
  {\Large\bfseries\color{op-nearblack}}
  {\thesubsection}{1em}{}

% ── Header / Footer ──
\usepackage{fancyhdr}
\pagestyle{fancy}
\fancyhf{}
\renewcommand{\headrulewidth}{0.4mm}
\renewcommand{\headrule}{\hbox to\headwidth{\color{op-rosa}\leaders\hrule height 0.4mm\hfill}}
\fancyhead[L]{\textcolor{op-textsec}{\footnotesize\bfseries ÓPERA PRIMA --- Plan de Entregas}}
\fancyhead[R]{\textcolor{op-textsec}{\footnotesize 4 de junio de 2026 --- Versión 2.0}}
\fancyfoot[C]{\textcolor{op-textmuted}{\footnotesize Página \thepage}}

% ── Links ──
\usepackage{hyperref}
\hypersetup{
  colorlinks=true,
  linkcolor=op-rosa,
  urlcolor=op-rosa,
}

% ── Lists ──
\usepackage{enumitem}
\setlist[itemize]{label=\textcolor{op-rosa}{\textbullet}, leftmargin=1.5em}

% ── Tables ──
\newcolumntype{L}[1]{>{\raggedright\arraybackslash}p{#1}}
\newcolumntype{C}[1]{>{\centering\arraybackslash}p{#1}}

% ── Misc ──
\usepackage{microtype}

% =============================================================================
\begin{document}

% ─── PORTADA ───
\thispagestyle{empty}

\begin{center}
  \vspace*{2.5cm}

  {\footnotesize\bfseries\color{op-rosa}
  \MakeUppercase{Propuesta de trabajo}}

  \vspace{0.6cm}

  {\Huge\bfseries\color{op-nearblack}
  Plan de Entregas\\[0.3em]
  Ópera Prima}

  \vspace{0.4cm}

  {\color{op-rosa}\rule{0.35\textwidth}{0.5mm}}

  \vspace{0.4cm}

  {\large\color{op-textsec}
  Plataforma digital para artistas emergentes}

  \vspace{2cm}

  \begin{tcolorbox}[op-box, width=0.6\textwidth, center]
    \begin{tabular}{>{\color{op-textsec}\bfseries}r @{\hspace{0.8em}} l}
      Cliente: & Ángela María Rodríguez \\
      Proyecto: & Plataforma Ópera Prima \\
      Duración: & 4 Jun $\rightarrow$ 4 Jul (30 días / 2 sprints) \\
      Fecha: & 4 de junio de 2026 \\
      Versión: & 2.0 \\
    \end{tabular}
  \end{tcolorbox}

  \vspace{1.5cm}

  \begin{tcolorbox}[op-box, width=0.6\textwidth, center]
    \small\centering
    {\bfseries\color{op-nearblack} Equipo de trabajo}\\[4pt]
    \begin{tabular}{>{\color{op-textsec}}l l}
      \textcolor{op-rosa}{$\blacksquare$} & \textbf{Santiago Montaño} --- Frontend \\
      \textcolor{op-verde}{$\blacksquare$} & \textbf{Camilo Sotelo} --- Backend \\
      \textcolor{op-lavanda}{$\blacksquare$} & \textbf{Alexander Paz} --- Coordinación \\
    \end{tabular}
  \end{tcolorbox}

  \vspace{2.5cm}

  {\color{op-verde}\rule{0.15\textwidth}{0.3mm}}
  \hspace{1em}
  {\color{op-rosa}\rule{0.3\textwidth}{0.5mm}}
  \hspace{1em}
  {\color{op-lavanda}\rule{0.15\textwidth}{0.3mm}}

  \vfill
\end{center}

\newpage

% ─── RESUMEN EJECUTIVO ───
\section{Resumen ejecutivo}

El desarrollo de la plataforma \textbf{Ópera Prima} se divide en \textbf{2 sprints} de 15 días cada uno:

\medskip

\begin{tcolorbox}[op-mvp]
  \centering
  {\Large\bfseries\color{op-verde} Sprint 1: MVP (días 1--15)}\\[2pt]
  {\color{op-textsec} 4 -- 18 de junio --- Santiago y Camilo trabajan en paralelo cada uno en su campo}\\[4pt]
  {\bfseries\color{op-nearblack} Objetivo:} Versión funcional de la plataforma con frontend y backend operativos.
\end{tcolorbox}

\medskip

\begin{tcolorbox}[op-aviso]
  \centering
  {\Large\bfseries\color{op-rosa-dark} Sprint 2: Correcciones y ajustes (días 16--30)}\\[2pt]
  {\color{op-textsec} 19 de junio -- 4 de julio --- Ambos trabajan juntos en pulir y desplegar}\\[4pt]
  {\bfseries\color{op-nearblack} Objetivo:} Integración, QA, assets finales y deploy a producción.
\end{tcolorbox}

\medskip

\begin{tcolorbox}[op-box]
  \begin{minipage}[c]{0.12\textwidth}
    \centering
    \Huge\bfseries\color{op-verde} 15
  \end{minipage}
  \hfill
  \begin{minipage}[c]{0.82\textwidth}
    {\large\bfseries\color{op-nearblack} Días para el MVP --- 4 al 18 de junio}\\[2pt]
    {\color{op-textsec} Santiago entrega correcciones UI + frontend de funcionalidades. Camilo entrega backend completo con APIs.}
  \end{minipage}
\end{tcolorbox}

\medskip

\begin{tcolorbox}[op-box]
  \begin{minipage}[c]{0.12\textwidth}
    \centering
    \Huge\bfseries\color{op-rosa} 15
  \end{minipage}
  \hfill
  \begin{minipage}[c]{0.82\textwidth}
    {\large\bfseries\color{op-nearblack} Días para correcciones y ajustes --- 19 de junio al 4 de julio}\\[2pt]
    {\color{op-textsec} Integración frontend-backend, QA completo, carga de assets y deploy a producción.}
  \end{minipage}
\end{tcolorbox}

\bigskip

\subsection*{Estructura general}

\begin{table}[h]
  \centering
  \small
  \renewcommand{\arraystretch}{1.3}
  \begin{tabular}{L{2cm} C{2.5cm} C{0.8cm} L{2.2cm} L{5.5cm}}
    \toprule
    \rowcolor{op-nearblack!90}
    {\color{op-white}\bfseries Semana} &
    {\color{op-white}\bfseries Fechas} &
    {\color{op-white}\bfseries Días} &
    {\color{op-white}\bfseries Sprint} &
    {\color{op-white}\bfseries Enfoque} \\
    \midrule
    \textbf{Semana 1} & 4 -- 10 Jun & 5 & \multirow{2}{*}{\textbf{MVP}} & Santiago: Correcciones UI + páginas nuevas \\
    \cline{1-3}\cline{5-5}
    \rowcolor{op-surface}
    \textbf{Semana 2} & 11 -- 17 Jun & 5 & & Santiago: Funcionalidades frontend + perfiles \\
    \midrule
    \textbf{Semana 3} & 18 -- 24 Jun & 5 & \multirow{2}{*}{\textbf{Ajustes}} & Ambos: Backend + correcciones \\
    \cline{1-3}\cline{5-5}
    \rowcolor{op-surface}
    \textbf{Semana 4} & 25 Jun -- 4 Jul & 5 & & Ambos: Integración + QA + Deploy \\
    \midrule
    \rowcolor{op-rosa!10}
    \textbf{Total} & \textbf{4 semanas} & \textbf{20} & \textbf{2 sprints} & \textbf{Plataforma completa operativa} \\
    \bottomrule
  \end{tabular}
\end{table}

\bigskip

\begin{tcolorbox}[op-mvp]
  \centering\small\bfseries\color{op-verde}
  \MakeUppercase{MVP listo el 18 de junio:} Correcciones UI aplicadas, frontend de funcionalidades clave montado, backend con APIs listas para integrar.
\end{tcolorbox}

\newpage

% ─── SEMANA 1 ───
\begin{tcolorbox}[op-fase]
  \makebox[0pt][l]{\hspace*{-5mm}{\color{op-rosa}\rule{2mm}{1.2em}}}%
  \hspace{3mm} Semana 1: MVP --- Frontend + Backend en paralelo
\end{tcolorbox}

\medskip
{\small\color{op-textsec} \textbf{4 -- 10 de junio de 2026} \hfill \textbf{Santiago (Front) + Camilo (Back) trabajan simultáneamente}}

\medskip

\textbf{Santiago} aplica las correcciones UI mientras \textbf{Camilo} configura la infraestructura backend. No se esperan el uno al otro.

\vspace{0.2cm}

\begin{tcolorbox}[op-aviso]
  \small\bfseries\color{op-rosa-dark}
  El cliente debe entregar: título del Evento Networking \#1 durante esta semana. La imagen 1920×1080 del hero se entrega en Semana 4.
\end{tcolorbox}

\vspace{0.3cm}

\begin{minipage}[t]{0.48\textwidth}
  \vspace{0pt}
  \subsubsection*{Santiago --- Frontend}
  \begin{table}[h]
    \centering
    \footnotesize
    \renewcommand{\arraystretch}{1.15}
    \begin{tabular}{c L{2.5cm} C{0.5cm}}
      \toprule
      \rowcolor{op-rosa!80}
      {\color{op-white}\bfseries \#} & {\color{op-white}\bfseries Tarea} & {\color{op-white}\bfseries D} \\
      \midrule
      1.1 & Términos y Condiciones & 1 \\
      \rowcolor{op-surface}
      1.2 & Eliminar grid overlay & 0.5 \\
      1.3 & Hero /sobre -- textos & 0.5 \\
      \rowcolor{op-surface}
      1.4 & Hero /sobre -- fondo & 0.5 \\
      1.5 & Sección Servicios & 2 \\
      \rowcolor{op-surface}
      1.6 & Quitar iconos Valores & 0.5 \\
      1.7 & Cards Equipo clickeables & 1 \\
      \rowcolor{op-surface}
      1.8 & Hero /mentorias & 1 \\
      1.9 & Paso a paso 6 pasos & 1 \\
      \rowcolor{op-surface}
      1.10 & Texto rotativo & 1 \\
      1.11 & Textos mentores & 0.5 \\
      \rowcolor{op-surface}
      1.12 & Rediseño /eventos & 1 \\
      1.13 & Evento Networking \#1 & 1 \\
      \rowcolor{op-surface}
      1.14 & Instagram & 0.5 \\
      1.15 & Google Drive & 1 \\
      \bottomrule
    \end{tabular}
  \end{table}
\end{minipage}
\hfill
\begin{minipage}[t]{0.48\textwidth}
  \vspace{0pt}
  \subsubsection*{Camilo --- Backend}
  \begin{table}[h]
    \centering
    \footnotesize
    \renewcommand{\arraystretch}{1.15}
    \begin{tabular}{c L{2.5cm} C{0.5cm}}
      \toprule
      \rowcolor{op-verde!80}
      {\color{op-white}\bfseries \#} & {\color{op-white}\bfseries Tarea} & {\color{op-white}\bfseries D} \\
      \midrule
      B.1 & Configurar Supabase & 1.5 \\
      \rowcolor{op-surface}
      B.2 & Schema de tablas + RLS & 1 \\
      B.3 & Migrar autenticación & 2 \\
      \rowcolor{op-surface}
      B.4 & API de perfiles & 1.5 \\
      B.5 & API de Oportunidades & 1.5 \\
      \rowcolor{op-surface}
      B.6 & API de Calendario & 1.5 \\
      B.7 & API de reserva mentorías & 1.5 \\
      \rowcolor{op-surface}
      B.8 & API de registro eventos & 1 \\
      B.9 & API de newsletter & 0.5 \\
      \rowcolor{op-surface}
      B.10 & Dominio & 1 \\
      \bottomrule
    \end{tabular}
  \end{table}
\end{minipage}

\vspace{0.2cm}
{\footnotesize\color{op-textmuted} \textbf{Total S1:} \textasciitilde12 días frontend + \textasciitilde12 días backend --- tareas altamente paralelizables.}

\newpage

% ─── SEMANA 2 ───
\begin{tcolorbox}[op-fase]
  \makebox[0pt][l]{\hspace*{-5mm}{\color{op-verde}\rule{2mm}{1.2em}}}%
  \hspace{3mm} Semana 2: MVP --- Funcionalidades Frontend + Backend
\end{tcolorbox}

\medskip
{\small\color{op-textsec} \textbf{11 -- 17 de junio de 2026} \hfill \textbf{Santiago completa frontend -- Camilo termina APIs}}

\medskip

\textbf{Santiago} construye las páginas funcionales. \textbf{Camilo} termina las APIs que faltan. Al final de la semana se tiene el MVP completo.

\vspace{0.3cm}

\begin{minipage}[t]{0.48\textwidth}
  \vspace{0pt}
  \subsubsection*{Santiago --- Frontend}
  \begin{table}[h]
    \centering
    \footnotesize
    \renewcommand{\arraystretch}{1.15}
    \begin{tabular}{c L{2.5cm} C{0.5cm}}
      \toprule
      \rowcolor{op-rosa!80}
      {\color{op-white}\bfseries \#} & {\color{op-white}\bfseries Tarea} & {\color{op-white}\bfseries D} \\
      \midrule
      2.1 & Sistema de perfiles & 2 \\
      \rowcolor{op-surface}
      2.2 & Conectar cards a perfiles & 1 \\
      2.3 & Tablero Oportunidades & 2 \\
      \rowcolor{op-surface}
      2.4 & Calendario Comunitario & 2 \\
      2.5 & Formulario reserva mentorías & 1.5 \\
      \rowcolor{op-surface}
      2.6 & Formulario registro eventos & 1.5 \\
      2.7 & Newsletter & 1 \\
      \rowcolor{op-surface}
      2.8 & QA responsive & 1 \\
      \bottomrule
    \end{tabular}
  \end{table}
\end{minipage}
\hfill
\begin{minipage}[t]{0.48\textwidth}
  \vspace{0pt}
  \subsubsection*{Camilo --- Backend}
  \begin{table}[h]
    \centering
    \footnotesize
    \renewcommand{\arraystretch}{1.15}
    \begin{tabular}{c L{2.5cm} C{0.5cm}}
      \toprule
      \rowcolor{op-verde!80}
      {\color{op-white}\bfseries \#} & {\color{op-white}\bfseries Tarea} & {\color{op-white}\bfseries D} \\
      \midrule
      B.11 & Refinar APIs existentes & 2 \\
      \rowcolor{op-surface}
      B.12 & Tests de integración & 1.5 \\
      B.13 & Documentación técnica & 1 \\
      \rowcolor{op-surface}
      B.14 & Ajustes de seguridad & 1 \\
      B.15 & Preparar deploy backend & 1 \\
      \rowcolor{op-surface}
      \multicolumn{3}{l}{\footnotesize\itshape\color{op-textsec} APIs completadas en Semana 1.} \\
      \bottomrule
    \end{tabular}
  \end{table}
\end{minipage}

\vspace{0.3cm}

\begin{tcolorbox}[op-mvp]
  \centering\small\bfseries\color{op-verde}
  \MakeUppercase{Fin del Sprint 1 (18 de junio): MVP entregado.} Santiago: todas las páginas montadas. Camilo: backend operativo con APIs listas. El producto es funcional --- siguiente paso: integrar.
\end{tcolorbox}

\newpage

% ─── SEMANA 3 ───
\begin{tcolorbox}[op-fase]
  \makebox[0pt][l]{\hspace*{-5mm}{\color{op-lavanda}\rule{2mm}{1.2em}}}%
  \hspace{3mm} Semana 3: Sprint 2 --- Correcciones y ajustes
\end{tcolorbox}

\medskip
{\small\color{op-textsec} \textbf{18 -- 24 de junio de 2026} \hfill \textbf{Santiago + Camilo --- Integración y pruebas}}

\medskip

Comienza el Sprint 2. Se integra frontend con backend y se corrigen todos los detalles.

\vspace{0.3cm}

\begin{table}[h]
  \centering
  \footnotesize
  \renewcommand{\arraystretch}{1.2}
  \begin{tabular}{c L{2.8cm} L{7cm} C{0.6cm} L{1.5cm}}
    \toprule
    \rowcolor{op-nearblack!90}
    {\color{op-white}\bfseries \#} &
    {\color{op-white}\bfseries Tarea} &
    {\color{op-white}\bfseries Descripción} &
    {\color{op-white}\bfseries Días} &
    {\color{op-white}\bfseries Resp.} \\
    \midrule
    3.1 & Integración parcial & Conectar primeras APIs con frontend. Auth real, perfiles vivos. & 2 & S + C \\
    \rowcolor{op-surface}
    3.2 & QA de flujos críticos & Probar registro, login, edición de perfil, mentorías. & 1.5 & Santiago \\
    3.3 & Imagen hero & Cargar imagen 1920×1080 del cliente en /sobre. & 0.5 & Santiago \\
    \rowcolor{op-surface}
    3.4 & Assets Google Drive & Integrar imágenes y fotos de mentores. & 1 & Santiago \\
    3.5 & Correcciones rápidas & Bugs y ajustes detectados en QA. & 2 & S + C \\
    \rowcolor{op-surface}
    3.6 & Ajustes responsive & Revisión en móvil, tablet y desktop. & 1 & Santiago \\
    \bottomrule
  \end{tabular}
\end{table}

\vspace{0.2cm}
{\footnotesize\color{op-textmuted} \textbf{Total estimado:} \textasciitilde8 días --- tareas en paralelo.}

\newpage

% ─── SEMANA 4 ───
\begin{tcolorbox}[op-fase]
  \makebox[0pt][l]{\hspace*{-5mm}{\color{op-nearblack}\rule{2mm}{1.2em}}}%
  \hspace{3mm} Semana 4: Sprint 2 --- Deploy y cierre
\end{tcolorbox}

\medskip
{\small\color{op-textsec} \textbf{25 de junio -- 4 de julio de 2026} \hfill \textbf{Santiago + Camilo --- Preparación final}}

\medskip

Últimos ajustes, deploy a producción y reunión de cierre con el cliente.

\vspace{0.3cm}

\begin{table}[h]
  \centering
  \footnotesize
  \renewcommand{\arraystretch}{1.2}
  \begin{tabular}{c L{2.8cm} L{7cm} C{0.6cm} L{1.5cm}}
    \toprule
    \rowcolor{op-nearblack!90}
    {\color{op-white}\bfseries \#} &
    {\color{op-white}\bfseries Tarea} &
    {\color{op-white}\bfseries Descripción} &
    {\color{op-white}\bfseries Días} &
    {\color{op-white}\bfseries Resp.} \\
    \midrule
    4.1 & Integración total & Conectar frontend completo con backend. Datos dinámicos. & 2 & S + C \\
    \rowcolor{op-surface}
    4.2 & QA completo & Todos los flujos: oportunidades, mentorías, eventos, newsletter. & 1.5 & Santiago \\
    4.3 & Correcciones finales & Bugs y ajustes de última hora. & 1.5 & S + C \\
    \rowcolor{op-surface}
    4.4 & Deploy Vercel & Variables de entorno, secrets, build y producción. & 1 & Camilo \\
    4.5 & Reunión final & Presentación de la plataforma. Validación y cierre. & --- & Ambos \\
    \bottomrule
  \end{tabular}
\end{table}

\vspace{0.3cm}

\begin{tcolorbox}[op-aviso]
  \centering\Large\bfseries\color{op-rosa-dark}
  4 de julio de 2026 --- Entrega final de la plataforma
\end{tcolorbox}

% ─── ENTREGABLES ───
\newpage
\section{Entregables por sprint}

\medskip

\begin{center}
\small
\begin{tikzpicture}
  \draw[op-verde, line width=0.5mm, rounded corners=0] (0,0) rectangle (16,2.5);
  \fill[op-verde] (0.5,1.25) circle (4mm);
  \node[color=op-white, font=\bfseries] at (0.5,1.25) {1};
  \node[anchor=west, color=op-nearblack, font=\bfseries] at (1.3,1.8) {Sprint 1 --- MVP (4 al 18 de junio)};
  \node[anchor=west, color=op-textsec, font=\small] at (1.3,1.2) {Santiago: Correcciones UI, /terminos, /sobre, /mentorias, /eventos, perfiles, tablero, calendario, formularios};
  \node[anchor=west, color=op-textsec, font=\small] at (1.3,0.7) {Camilo: Supabase, auth real, APIs de perfiles, oportunidades, calendario, mentorías, eventos, newsletter};
  \node[anchor=west, color=op-textsec, font=\small] at (1.3,0.2) {\textbf{Resultado:} Plataforma funcional --- frontend y backend listos para integrar};

  \draw[op-nearblack, dashed, line width=0.2mm] (-0.3,-0.5) -- (16.3,-0.5);

  \draw[op-rosa, line width=0.5mm, rounded corners=0] (0,-1.3) rectangle (16,-3.8);
  \fill[op-rosa] (0.5,-2.55) circle (4mm);
  \node[color=op-white, font=\bfseries] at (0.5,-2.55) {2};
  \node[anchor=west, color=op-nearblack, font=\bfseries] at (1.3,-1.6) {Sprint 2 --- Correcciones y ajustes (19 de junio al 4 de julio)};
  \node[anchor=west, color=op-textsec, font=\small] at (1.3,-2.3) {Integración frontend-backend, QA completo, imagen hero, assets, correcciones};
  \node[anchor=west, color=op-textsec, font=\small] at (1.3,-2.8) {Deploy a producción, reunión final con el cliente};
  \node[anchor=west, color=op-textsec, font=\small] at (1.3,-3.3) {\textbf{Resultado:} Plataforma completa en producción};

  \draw[op-nearblack, dashed, line width=0.2mm] (-0.3,-4.6) -- (16.3,-4.6);

  \draw[op-nearblack, line width=0.2mm, rounded corners=0] (4.5,-5.3) rectangle (11.5,-6.8);
  \node[color=op-rosa, font=\Huge\bfseries] at (8,-6.05) {4 de julio};
\end{tikzpicture}
\end{center}

% ─── DEPENDENCIAS ───
\newpage
\section{Dependencias del cliente}

\medskip

\begin{table}[h]
  \centering
  \small
  \renewcommand{\arraystretch}{1.2}
  \begin{tabular}{L{5cm} C{3cm} L{6cm}}
    \toprule
    \rowcolor{op-nearblack!90}
    {\color{op-white}\bfseries Elemento} &
    {\color{op-white}\bfseries Fecha tope} &
    {\color{op-white}\bfseries ¿Qué bloquea?} \\
    \midrule
    Imagen hero 1920×1080 &
    19 de junio &
    Hero /sobre se queda con fondo negro \\
    \rowcolor{op-surface}
    Assets en Google Drive &
    11 de junio &
    No se integran elementos visuales \\
    Definición del dominio &
    25 de junio &
    Deploy a producción \\
    \rowcolor{op-surface}
    Título Evento Networking \#1 &
    11 de junio &
    Card se publica con ``Por definir'' \\
    Disponibilidad reuniones &
    18 jun y 4 jul &
    Validación del MVP y cierre \\
    \rowcolor{op-surface}
    Datos reales de mentores &
    11 de junio &
    Perfiles con placeholder \\
    \bottomrule
  \end{tabular}
\end{table}

% ─── NOTAS ───
\section{Notas importantes}

\medskip

\begin{enumerate}
  \item \textbf{Días calendario, no hábiles.} El plazo es del 4 de junio al 4 de julio.
  \item \textbf{Trabajo en paralelo desde el día 1.} Santiago y Camilo arrancan juntos pero cada uno en su campo. No se bloquean mutuamente.
  \item \textbf{Supabase} tiene plan gratuito que cubre las necesidades iniciales.
  \item \textbf{Los pagos y facturación} están a cargo de Alexander Paz y no forman parte de este cronograma.
  \item Se asume \textbf{respuesta del cliente en máximo 48 horas} para dudas y aprobaciones.
  \item Los elementos visuales siguen la línea editorial del manual de diseño (\texttt{DESIGN.md}).
\end{enumerate}

\vspace{1.5cm}

\begin{center}
  {\color{op-verde}\rule{0.2\textwidth}{0.3mm}}
  \hspace{1.5em}
  {\color{op-textmuted}\footnotesize Documento preparado para revisión del cliente --- Versión 2.0}
  \hspace{1.5em}
  {\color{op-rosa}\rule{0.2\textwidth}{0.3mm}}
\end{center}

\end{document}
